% hole_top_mode.tex -- round 413, one lemma:
% the top OP of the hole measure, and why codimension-1
% contractivity of T0 on that mode's complement is not P1.
% TOP_MODE_REFUTED.
% Builds with pdflatex.
% Companion machine check: verify_hole_top_mode.py.
% Discovery: experiments/tfpt-discovery/hole_top_mode_probe.py.
% NO RH CLAIM.  Research documentation, not a theorem of RH.
\documentclass[11pt]{article}

\usepackage[a4paper,margin=2.7cm]{geometry}
\usepackage{amsmath,amssymb,amsthm}
\usepackage{microtype}
\usepackage{booktabs}
\usepackage{xcolor}
\usepackage[colorlinks=true,linkcolor=blue!50!black,urlcolor=blue!50!black]{hyperref}

\newtheorem{lemma}{Lemma}
\newtheorem{prop}{Proposition}
\theoremstyle{remark}
\newtheorem{remark}{Remark}

\newcommand{\Tfrak}{\mathfrak{T}}

\title{\bfseries The hole top mode\\[2mm]
\large $v_{\mathrm{top}}\propto 1/(\sqrt{\omega}\,P_Y')$ is the
top hole OP; it is not the excess of $\Tfrak_0$}
\author{Stefan Hamann}
\date{August 29, 2026}

\begin{document}
\maketitle

\begin{abstract}\noindent
After r407--r411 the remaining P1-shaped statement is
``at most one negative direction''.
The candidate for that direction is the top orthogonal
polynomial of the hole measure $(Y,u^\vee)$:
$v_{\mathrm{top},i}\propto 1/(\sqrt{u^\vee(y_i)}\,P_Y'(y_i))$,
source-pure from $Y$, $u^\vee$, $P_Y'$
(no eigenvector, no SVD, no target).
The formula is a theorem over $\mathbb{Q}$.
The Hole-Top-Mode test
$Q_{\mathrm{top}}(I-\Tfrak_0^*\Tfrak_0)Q_{\mathrm{top}}\succeq 0$
\textbf{fails on flagship MAIN}
($\lambda_{\min}=-0.0476$, leftover $\sigma=1.024$ on $w9$).
Quadrature dominance already fails for $p\equiv 1$.
One clean MAIN violation closes the route.
No RH claim.
\end{abstract}

\medskip\noindent
\textbf{Status.}\enspace
Lemmas~\ref{lem:lagrange}--\ref{lem:formula} are proved over
$\mathbb{Q}$.
Proposition~\ref{prop:minmax} is linear algebra
(the implication is not in doubt; its hypothesis is false).
HTM on MAIN, QD, and ``dead $\chi$ hold'' are
\textbf{refuted}.
Ausgang \texttt{TOP\_MODE\_REFUTED}.
No $L^*$ claim.  No $R^\dagger$ claim.  No RH claim.

\section{The formula}

Let $Y=\{y_1,\dots,y_q\}$ be simple, $\omega_i>0$,
$P_Y(x)=\prod_k(x-y_k)$.

\begin{lemma}[Lagrange identity]\label{lem:lagrange}
For every polynomial $p$ of degree at most $q-2$,
\[
  \sum_{i=1}^q \frac{p(y_i)}{P_Y'(y_i)}\;=\;0.
\]
On four rational atoms this is exact in
$\mathbb{Q}$ for the monomial, Chebyshev, and Newton
bases.
\end{lemma}

The residue of $p/P_Y$ at infinity vanishes in that
degree range, and the partial-fraction coefficient of
$1/(x-y_i)$ is $p(y_i)/P_Y'(y_i)$.

\begin{lemma}[Top hole OP]\label{lem:formula}
The samples $\pi_{q-1}^Y(y_i)=c/(\omega_i P_Y'(y_i))$
are orthogonal to every polynomial of degree
$\le q-2$ in the discrete inner product
$\langle f,g\rangle_\omega=\sum_i\omega_i f(y_i)g(y_i)$.
Gram--Schmidt in three bases yields the same samples
up to scale.
In the $\sqrt{\omega}$-coordinate,
\[
  v_{\mathrm{top},i}
  \;\propto\;
  \frac{1}{\sqrt{\omega_i}\,P_Y'(y_i)}.
\]
\end{lemma}

The constructor uses only $Y$ and $u^\vee=\omega$.
No eigendecomposition of $C$, no SVD of $\Tfrak_0$,
no target matrix.

\section{The primary test}

$\Tfrak_0$ is the min-norm interpolant $Y\to X$ in
$\mathcal{P}_{<d_0}$ (r409).
Write $Q=I-vv^T$ for unit $v=v_{\mathrm{top}}$.

\begin{prop}[Min-max]\label{prop:minmax}
If $\|\Tfrak_0 z\|_X\le\|z\|_Y$ for all $z\perp v_{\mathrm{top}}$,
then $\Tfrak_0$ has at most one singular value $>1$.
By the r407/r411 dictionary $\Tfrak_0^*\Tfrak_0=C^{-1}$
this is $\lambda_2(C)\ge 1$, which is P1.
\end{prop}

The hypothesis is the projected Loewner condition
$Q(I-\Tfrak_0^*\Tfrak_0)Q\succeq 0$, equivalently
$\#\{\sigma(\Tfrak_0|_{v^\perp})>1\}=0$.

\paragraph{Flagship MAIN $w9$.}
$\operatorname{corr}(v_{\mathrm{top}},C\text{-ev}_0)=0.6766$,
better than the r410 Fourier face ($0.386$ vs $C$-ev$_0$,
$0.278$ vs $v_{\mathrm{top}}$).
Projecting $v_{\mathrm{top}}$ cuts
$\|\Tfrak_0\|=1.08014$ to $\sigma_{\max}(v^\perp)=1.02354$,
still $>1$: $\lambda_{\min}=-0.04763$, $n_{>1}=1$.
The Nyquist face is real and is not the excess
(r411: the excess \emph{is} $C$-ev$_0$).

\paragraph{Census.}
Core-$42$: P1 windows hold $4/28$, PD $14/14$ (trivial:
$\|\Tfrak_0\|\le 1$ already).
EXT P1 $\{119,35,109,69\}$ all break.
Living $\chi$: P1 $1/37$, PD $41/41$.
Dead $\chi$ hold $4/6$; the two dead-P1 windows
$\chi_3$-$19$ and $\chi_3$-$39$ break.
Death is not an HTM-edge.

\section{Quadrature dominance}

Every $z\perp v_{\mathrm{top}}$ is the $\sqrt{\omega}$-sample
of a unique $p\in\mathcal{P}_{\le q-2}$.
QD asks $\sum_X u^\vee p^2\le\sum_Y u^\vee p^2$.
Min-norm gives $\|\Tfrak_0 z\|\le\|p\|_X$, so QD implies HTM.

On $w9$, already $p\equiv 1$ fails:
$\sum_X u^\vee\big/\sum_Y u^\vee=1.35713$.
The min-norm gap \emph{does} save this direction
($\|\Tfrak_0 v_1\|=0.359<1$) and does \emph{not}
save the leftover excess ($\sigma_{\max}(v^\perp)>1$).
QD is strictly stronger than HTM, and both fail.

\section{Kills}

Permutation (same $Y$) and scramble break the primary
test with $n_{>1}=20$ and $\lambda_{\min}\in\{-54,-217\}$.
Dropping the signs of $P_Y'$ leaves
$\sigma_{\max}\approx\|\Tfrak_0\|$ (the Nyquist face dies).
A random unit vector does the same.
Substituting raw $u$ for $u^\vee$ drops the $C$-ev$_0$
correlation to $0.126$ and leaves $n_{>1}=1$.

\section{What this closes}

The formula is a theorem, and $v_{\mathrm{top}}$ is the
canonical top hole mode.
It is not the one negative direction of P1.
The reviewer order applies: one clean MAIN violation
closes the hole-top-mode route.
P1 remains $\iff\lambda_2(C)\ge 1$; that RHS is not this
vector.
No RH claim.

\end{document}
